% Nouveau EGA V (cote 25): archive cote 25, site page 120 (PDF page 121).
% typescript. Draft compared: none.
% AI-transcribed by Claude Opus 5.5 (claude-opus-5-5, adaptive thinking, effort high), Message Batches API, batch ega25b (msgbatch_01A7n5YuXubUPNK1zsepi8RG), 2026-11-13; not yet checked by eye.
\scanpage{25/120}
\mspage{\unsure{27}}

Z$_s$ est de codimension $\geq 2$ dans X$_s$ ; \var{\xout{\illegible{2}\unsure{(comme dans les cas envisa}}}{}{typewriter strike-out}{ega25_25_120-v0001}
\var{\xout{familles}}{}{typewriter strike-out}{ega25_25_120-v0002} \var{\inserted{syst lin}}{syst lin}{insertion}{ega25_25_120-v0003} de diviseurs par.par P qui so\var{\inserted{\unsure{ie}}}{\unsure{ie}}{insertion}{ega25_25_120-v0004}nt
\uline{sans composante fixe}, \var{\xout{tandis que $\emptyset$ est l'ensemble de}}{}{typewriter strike-out}{ega25_25_120-v0005} D'autre part, si
S=Spec(k), et si X est normal, alors \var{\inserted{pour}}{pour}{insertion}{ega25_25_120-v0006} \uline{toute} application rationelle
f:X$\to$P, l'ensemble de définition U(f) vérifie codim ((X-U(f),X)$\geq$2
(II 7), donc dans ce cas $\emptyset$ est formé de l'ensemble de \uline{toutes} les
applications rationelles de X dans P.

La démonstration de a) a déjà été donnée. Pour prouver b), notons
que la formation de $i_*(\underline{L}_U)$ commute avec toute extension plate de la
\var{\xout{pour le\unsure{x} il faut}}{}{typewriter strike-out}{ega25_25_120-v0007} \var{\inserted[typed]{la cond. envisagée est inv. par}}{la cond. envisagée est inv. par}{insertion}{ega25_25_120-v0008}
base S' \quad S (réf)\var{\inserted{\unsure{du moins si $U\to X$ est quasi-compact, cas auquel on se ramène sans difficultés}}}{\unsure{du moins si $U\to X$ est quasi-compact, cas auquel on se ramène sans difficultés}}{insertion}{ega25_25_120-v0009}, donc \var{\xout{il suffit de trouver un morphisme}}{}{typewriter strike-out}{ega25_25_120-v0010} fid plat qu cpc
\var{\xout{S' $\quad$ S tel que la conclusion soit vraie après}}{}{typewriter strike-out}{ega25_25_120-v0011} changement de base. On
prendra S' = $\hand{\check{}}\mathrm{P}$ , et on note que l'hypothèse $i'_*(\underline{L}'_U)$ inversible ne
change pas en remplaçant $\underline{L}'_U$ par $\underline{L}'_U\otimes_{S'}M'$, où M' est un Module inversible
sur S' (car $i'_*(\underline{L}'_U\otimes_{S'}M')\simeq i'_*(\underline{L}'_U)\otimes_{S'}M'$ ). On \unsure{prèndra|prendra} M'= $\underline{O}_{\check{P}}(1)$, de
sorte que la condition envisagée signifie aussi que $i'_*(\underline{N}')$ est un Module
inversible, où $\underline{N}'= (f_U\times_S id_P)^*(\underline{O}_{P\times\check{P}}(1,1))$ . Or $\underline{O}$ (1,1) est précisément
le Module inversible défini par le diviseur canonique $\underline{H}$ sur $P\times_S\check{P}$ , de sor%
te que $\underline{N}$ n'est autre que le Module inversible défini par \qquad D' =
$(f_U\times id_P)$ $(\underline{H})$. \var{\deleted{\var{\xout{\illegible{1}}}{}{typewriter strike-out}{ega25_25_120-v0013}}}{}{deletion}{ega25_25_120-v0012} Lorsque f provient de D$\in\emptyset$ , \var{\xout{ce dernier}}{}{typewriter strike-out}{ega25_25_120-v0014}
D' n'est autre que D$|$U$\times_S\check{P}$ , \var{\xout{et il résulte du lemme}}{}{typewriter strike-out}{ega25_25_120-v0015} donc $\underline{N}' = \underline{N}|U\times_S\check{P}$ ,
où $\underline{N}$ est le Module inversible sur $X\times_S\hand{\check{}}P$ défini par D, et il résulte
du lemme \var{\inserted{3.2}}{3.2}{insertion}{ega25_25_120-v0016} plus haut \var{\xout{\illegible{1}}}{}{typewriter strike-out}{ega25_25_120-v0017} appliqué à $P\times_S\hand{\check{}}P\to\hand{\check{}}P$ qu'on aura
$i'_*(\underline{N}')\simeq\underline{N}$ , donc $i'_*(\underline{N}')$ est inversible. Inversement, si cette condi-%
tion est remplie prouvons que f provient d'un\var{\xout{e}}{}{typewriter strike-out}{ega25_25_120-v0018} D$\in\emptyset$ , ou ce qui revient
ma\var{\xout{n}}{}{typewriter strike-out}{ega25_25_120-v0019}nifestement au même, que le diviseur D' se prolonge en un diviseur \var{\deleted{de}}{}{deletion}{ega25_25_120-v0020}
\var{\deleted{Cartier}}{}{deletion}{ega25_25_120-v0021} relatif par r apport \var{\xout{de}}{}{typewriter strike-out}{ega25_25_120-v0022} à $\hand{\check{}}P$ sur X$\times_S$P . Il \var{\deleted{suffit}}{}{deletion}{ega25_25_120-v0023} revient au même
de dire qu'il se prolonge en un diviseur \var{\deleted{de Cartier}}{}{deletion}{ega25_25_120-v0024} D sur X$\times_S$P , car \hand{D}.

\var{\margin{\unsure{Rq compliquée} \unsure{Il faut expliquer} \illegible{2} \unsure{$f\times\mathrm{id}$} \unsure{si $p^{-1}(H)$} \unsure{est défini} \illegible{1} \unsure{diviseur rel.} \unsure{à la condition} \illegible{3} \unsure{diviseur relatif} \unsure{$X\times P$} \unsure{/P} \illegible{2}}}{\placedmargin{\unsure{Rq compliquée} \unsure{Il faut expliquer} \illegible{2} \unsure{$f\times\mathrm{id}$} \unsure{si $p^{-1}(H)$} \unsure{est défini} \illegible{1} \unsure{diviseur rel.} \unsure{à la condition} \illegible{3} \unsure{diviseur relatif} \unsure{$X\times P$} \unsure{/P} \illegible{2}}}{marginal note placed}{ega25_25_120-v0025}

\mspage{120}
